\chapter{Spheres and Antipodal Maps}
\label{ch:viii04-spheres-antipodal}

Spheres are the first family on which homotopy and degree interact systematically.
The involution \(x\mapsto -x\) exposes a parity phenomenon, while the same
geometry leads to Borsuk--Ulam obstructions, projective quotients, and the cell
structures used in the next chapters.

\section{Spheres, disks, and hemispheres}
Write
\[
S^n=\{x\in\mathbb R^{n+1}:\|x\|=1\},\qquad
D^{n+1}=\{x\in\mathbb R^{n+1}:\|x\|\le1\}.
\]
Then \(\partial D^{n+1}=S^n\).  The upper and lower closed hemispheres of
\(S^n\) are copies of \(D^n\) meeting in the equatorial \(S^{n-1}\).

\begin{proposition}[Two-disk description]\label{prop:viii04-two-disks}
For \(n\ge1\), \(S^n\) is obtained by gluing two copies of \(D^n\) along their
boundary \(S^{n-1}\) by the identity.
\end{proposition}

\section{The antipodal map}
\begin{definition}[Antipodal map]\label{def:viii04-antipodal}
The antipodal map is
\[
a:S^n\to S^n,\qquad a(x)=-x.
\]
It is a fixed-point-free involution: \(a^2=\operatorname{id}\).
\end{definition}

\begin{theorem}[Antipodal degree]\label{thm:viii04-antipodal-degree}
With the standard orientation,
\[
\deg(a)=(-1)^{n+1}.
\]
\end{theorem}

\begin{corollary}[Parity obstruction]\label{cor:viii04-parity}
If \(n\) is even, \(a\) is not homotopic to the identity map of \(S^n\).
\end{corollary}

\section{Odd and even maps}
\begin{definition}[Odd and even]\label{def:viii04-odd-even}
A map \(f:S^n\to\mathbb R^m\) is odd if \(f(-x)=-f(x)\), and even if
\(f(-x)=f(x)\).
\end{definition}

\begin{theorem}[Odd self-map theorem]\label{thm:viii04-odd-degree}
Every continuous odd map \(f:S^n\to S^n\) has odd degree.
\end{theorem}

\begin{corollary}\label{cor:viii04-odd-nonnull}
An odd self-map of \(S^n\) is not null-homotopic.
\end{corollary}

\section{Borsuk--Ulam}
\begin{theorem}[Borsuk--Ulam]\label{thm:viii04-borsuk-ulam}
For every continuous \(f:S^n\to\mathbb R^n\), there exists \(x\in S^n\) such
that
\[
f(x)=f(-x).
\]
\end{theorem}

\begin{theorem}[Equivalent no-odd-map form]\label{thm:viii04-no-odd-lower}
There is no continuous odd map \(S^n\to S^{n-1}\).
\end{theorem}

If a map \(f:S^n\to\mathbb R^n\) separated every antipodal pair, then
\[
g(x)=\frac{f(x)-f(-x)}{\|f(x)-f(-x)\|}
\]
would be an odd map \(S^n\to S^{n-1}\).  This proves the equivalence of the
two formulations.

\section{The one-dimensional case}
For \(f:S^1\to\mathbb R\), put \(h(x)=f(x)-f(-x)\).  Then
\(h(-x)=-h(x)\), so the intermediate value theorem on a semicircle forces a
zero of \(h\).  Thus some antipodal pair has the same value.

\section{No retraction}
\begin{theorem}[No boundary retraction]\label{thm:viii04-no-retraction}
There is no retraction \(D^{n+1}\to S^n\).
\end{theorem}
For \(n\ge1\), degree proves this: if \(r\) were a retraction and
\(i:S^n\hookrightarrow D^{n+1}\), then \(ri=\operatorname{id}\) has degree
\(1\), while \(i\) is null-homotopic in the disk and hence \(ri\) would have
degree \(0\).

\section{Real projective space}
\begin{definition}[Real projective space]\label{def:viii04-rpn}
\[
\mathbb RP^n=S^n/(x\sim -x).
\]
Equivalently, \(\mathbb RP^n\) is the set of one-dimensional linear subspaces
of \(\mathbb R^{n+1}\).
\end{definition}

Let \(q:S^n\to\mathbb RP^n\) be the quotient map.

\begin{proposition}[Even-map factorization]\label{prop:viii04-factor}
If \(f:S^n\to Y\) is continuous and even, then it factors uniquely as
\[
S^n\xrightarrow{q}\mathbb RP^n\xrightarrow{\bar f}Y.
\]
\end{proposition}

\section{Low-dimensional projective spaces}
\[
\mathbb RP^0\cong\{*\},\qquad \mathbb RP^1\cong S^1.
\]
Under a standard identification \(\mathbb RP^1\cong S^1\), the quotient map is
the degree-two circle map \(z\mapsto z^2\).

\section{Projective filtration}
There are natural inclusions
\[
\mathbb RP^0\subset\mathbb RP^1\subset\cdots\subset\mathbb RP^n.
\]
Moreover,
\[
\mathbb RP^k\setminus\mathbb RP^{k-1}\cong\mathbb R^k,
\]
so \(\mathbb RP^n\) has one open cell in each dimension \(0,\dots,n\).

\section{One-point compactification}
Stereographic projection gives
\[
S^n\setminus\{N\}\cong\mathbb R^n,
\]
so \(S^n\) is the one-point compactification of \(\mathbb R^n\).

\section{Suspension}
The sphere also satisfies
\[
S^{n+1}\cong\Sigma S^n.
\]
This viewpoint connects sphere geometry to cones and mapping cones.

\section{Bridge to cell attachment}
The projective filtration can be written schematically as
\[
\mathbb RP^n=\mathbb RP^{n-1}\cup e^n.
\]
VIII/05 formalizes this operation as attachment of an \(n\)-disk along its
boundary.

\begin{warning}[Quotient versus subspace]\label{warn:viii04-quotient}
The antipodal quotient identifies points; it is not an inclusion.  Retraction,
subspace, and quotient arguments must therefore be kept distinct.
\end{warning}


\section{Exercises}

\begin{exercise}\label{exr:viii04-01}
Compute \(a^2\) for \(a(x)=-x\).
\end{exercise}
\begin{hint}
Apply the map twice.

% VIII-PEDAGOGY-HINT-v1.1 exr:viii04-01 append
Compute the two successive ambient sign changes, and check that both intermediate and final vectors still lie on the unit sphere.
\end{hint}
\begin{solution}
\(a(a(x))=x\).
\end{solution}

\begin{exercise}\label{exr:viii04-02}
Does \(a\) have a fixed point on \(S^n\)?
\end{exercise}
\begin{hint}
Solve \(x=-x\).

% VIII-PEDAGOGY-HINT-v1.1 exr:viii04-02 append
Solve the fixed-point equation in the ambient real vector space, then impose the sphere's unit-norm condition.
\end{hint}
\begin{solution}
No; that would force \(x=0\).
\end{solution}

\begin{exercise}\label{exr:viii04-03}
Find \(\deg a\) on \(S^4\).
\end{exercise}
\begin{hint}
Use \((-1)^{n+1}\).

% VIII-PEDAGOGY-HINT-v1.1 exr:viii04-03 append
Use the determinant of \(-I\) in \(\mathbb R^5\); the exponent counts ambient coordinates, not the sphere dimension alone.
\end{hint}
\begin{solution}
\(-1\).
\end{solution}

\begin{exercise}\label{exr:viii04-04}
Find \(\deg a\) on \(S^5\).
\end{exercise}
\begin{hint}
Use parity.

% VIII-PEDAGOGY-HINT-v1.1 exr:viii04-04 append
There are six ambient coordinate sign reversals for \(S^5\). Relate their parity to the orientation sign.
\end{hint}
\begin{solution}
\(1\).
\end{solution}

\begin{exercise}\label{exr:viii04-05}
When does degree rule out \(a\simeq\operatorname{id}\)?
\end{exercise}
\begin{hint}
Compare their degrees.

% VIII-PEDAGOGY-HINT-v1.1 exr:viii04-05 append
Use homotopy invariance to obtain a necessary equality of degrees. Equality by itself is not being used here as a general classification theorem.
\end{hint}
\begin{solution}
For even \(n\).
\end{solution}

\begin{exercise}\label{exr:viii04-06}
State Borsuk--Ulam on \(S^2\).
\end{exercise}
\begin{hint}
The target is \(\mathbb R^2\).

% VIII-PEDAGOGY-HINT-v1.1 exr:viii04-06 append
Specify that the assertion concerns every continuous map and at least one pair \(x,-x\), with equality in both real coordinates simultaneously.
\end{hint}
\begin{solution}
Every continuous \(S^2\to\mathbb R^2\) identifies an antipodal pair.
\end{solution}

\begin{exercise}\label{exr:viii04-07}
Why would an odd \(S^n\to S^{n-1}\) violate Borsuk--Ulam?
\end{exercise}
\begin{hint}
Embed \(S^{n-1}\) in \(\mathbb R^n\).

% VIII-PEDAGOGY-HINT-v1.1 exr:viii04-07 append
Equal antipodal images combined with oddness would force the common image to be zero; test whether zero belongs to the unit-sphere target.
\end{hint}
\begin{solution}
Oddness prevents equal antipodal images.
\end{solution}

\begin{exercise}\label{exr:viii04-08}
Construct an odd map from an antipode-separating \(f:S^n\to\mathbb R^n\).
\end{exercise}
\begin{hint}
Normalize \(f(x)-f(-x)\).

% VIII-PEDAGOGY-HINT-v1.1 exr:viii04-08 append
First prove the denominator is nowhere zero from antipode separation. Then check continuity and that the normalized vector has unit norm.
\end{hint}
\begin{solution}
\(g(x)=\frac{f(x)-f(-x)}{\|f(x)-f(-x)\|}\).
\end{solution}

\begin{exercise}\label{exr:viii04-09}
Show that this \(g\) is odd.
\end{exercise}
\begin{hint}
Replace \(x\) by \(-x\).

% VIII-PEDAGOGY-HINT-v1.1 exr:viii04-09 append
Compute the numerator and its norm separately after replacing \(x\) by \(-x\); the norm is unchanged by sign reversal.
\end{hint}
\begin{solution}
Its numerator changes sign and denominator does not.
\end{solution}

\begin{exercise}\label{exr:viii04-10}
State Borsuk--Ulam for \(S^1\to\mathbb R\).
\end{exercise}
\begin{hint}
Set \(n=1\).

% VIII-PEDAGOGY-HINT-v1.1 exr:viii04-10 append
Set \(h(x)=f(x)-f(-x)\) and restrict to a semicircle joining antipodes. Its endpoint values allow an intermediate-value argument.
\end{hint}
\begin{solution}
Some antipodal pair has equal value.
\end{solution}

\begin{exercise}\label{exr:viii04-11}
Define \(\mathbb RP^n\) as a quotient.
\end{exercise}
\begin{hint}
Identify antipodes.

% VIII-PEDAGOGY-HINT-v1.1 exr:viii04-11 append
Describe both the equivalence relation and the quotient topology; naming the orbit set alone does not specify the topological space.
\end{hint}
\begin{solution}
\(\mathbb RP^n=S^n/(x\sim -x)\).
\end{solution}

\begin{exercise}\label{exr:viii04-12}
Give the linear-algebra description of \(\mathbb RP^n\).
\end{exercise}
\begin{hint}
Use lines through the origin.

% VIII-PEDAGOGY-HINT-v1.1 exr:viii04-12 append
Intersect a line through the origin with \(S^n\); determine the pair of unit representatives and compare it with one antipodal equivalence class.
\end{hint}
\begin{solution}
One-dimensional linear subspaces of \(\mathbb R^{n+1}\).
\end{solution}

\begin{exercise}\label{exr:viii04-13}
When does \(f:S^n\to Y\) factor through \(q\)?
\end{exercise}
\begin{hint}
Be constant on quotient fibers.

% VIII-PEDAGOGY-HINT-v1.1 exr:viii04-13 append
Define a candidate map on classes by \([x]\mapsto f(x)\). Separate well-definedness on each two-point fiber from continuity supplied by the quotient property.
\end{hint}
\begin{solution}
When \(f(x)=f(-x)\).
\end{solution}

\begin{exercise}\label{exr:viii04-14}
Identify \(\mathbb RP^0\).
\end{exercise}
\begin{hint}
There is one line in \(\mathbb R\).

% VIII-PEDAGOGY-HINT-v1.1 exr:viii04-14 append
Compare the two points of \(S^0\) with their single orbit under the antipodal involution.
\end{hint}
\begin{solution}
A point.
\end{solution}

\begin{exercise}\label{exr:viii04-15}
Identify \(\mathbb RP^1\).
\end{exercise}
\begin{hint}
Antipodal circle quotient.

% VIII-PEDAGOGY-HINT-v1.1 exr:viii04-15 append
Represent a line by an angle modulo \(\pi\) and double the angle to obtain a well-defined point on the circle.
\end{hint}
\begin{solution}
It is homeomorphic to \(S^1\).
\end{solution}

\begin{exercise}\label{exr:viii04-16}
What circle map models \(q:S^1\to\mathbb RP^1\cong S^1\)?
\end{exercise}
\begin{hint}
Two-to-one covering.

% VIII-PEDAGOGY-HINT-v1.1 exr:viii04-16 append
Seek a circle map unchanged when \(z\) is replaced by \(-z\); then check that its fibers contain exactly those two points.
\end{hint}
\begin{solution}
\(z\mapsto z^2\).
\end{solution}

\begin{exercise}\label{exr:viii04-17}
What is the degree of that map?
\end{exercise}
\begin{hint}
Use the power-map theorem.

% VIII-PEDAGOGY-HINT-v1.1 exr:viii04-17 append
Use the angle-doubling identification from the preceding exercise with its chosen orientation; an arbitrary orientation change of the target changes the degree's sign.
\end{hint}
\begin{solution}
\(2\).
\end{solution}

\begin{exercise}\label{exr:viii04-18}
Why is \(S^n\) a one-point compactification of \(\mathbb R^n\)?
\end{exercise}
\begin{hint}
Use stereographic projection.

% VIII-PEDAGOGY-HINT-v1.1 exr:viii04-18 append
After stereographic projection, verify that neighborhoods of the omitted pole correspond to complements of compact subsets of \(\mathbb R^n\).
\end{hint}
\begin{solution}
Removing one point gives \(\mathbb R^n\).
\end{solution}

\begin{exercise}\label{exr:viii04-19}
Express \(S^{n+1}\) as a suspension.
\end{exercise}
\begin{hint}
Suspend \(S^n\).

% VIII-PEDAGOGY-HINT-v1.1 exr:viii04-19 append
Use height \(t\in[-1,1]\) and the map \((x,t)\mapsto(\sqrt{1-t^2}x,t)\); examine what happens at the two extreme heights.
\end{hint}
\begin{solution}
\(S^{n+1}\cong\Sigma S^n\).
\end{solution}

\begin{exercise}\label{exr:viii04-20}
How many standard cells does \(\mathbb RP^n\) have?
\end{exercise}
\begin{hint}
One in each dimension \(0,\dots,n\).

% VIII-PEDAGOGY-HINT-v1.1 exr:viii04-20 append
Count the dimensions in the filtration \(\mathbb RP^0\subset\cdots\subset\mathbb RP^n\), remembering the zero-dimensional cell.
\end{hint}
\begin{solution}
\(n+1\).
\end{solution}

\begin{exercise}\label{exr:viii04-21}
What is the equator of \(S^n\)?
\end{exercise}
\begin{hint}
Set the last coordinate to zero.

% VIII-PEDAGOGY-HINT-v1.1 exr:viii04-21 append
With the last coordinate fixed to zero, the remaining coordinates still satisfy one unit-sphere equation in \(\mathbb R^n\).
\end{hint}
\begin{solution}
A copy of \(S^{n-1}\).
\end{solution}

\begin{exercise}\label{exr:viii04-22}
Why is a closed hemisphere a disk?
\end{exercise}
\begin{hint}
Project to the first \(n\) coordinates.

% VIII-PEDAGOGY-HINT-v1.1 exr:viii04-22 append
On the upper hemisphere recover the last coordinate as \(\sqrt{1-\|u\|^2}\). This constructs the inverse to projection, including the boundary.
\end{hint}
\begin{solution}
It is homeomorphic to \(D^n\).
\end{solution}

\begin{exercise}\label{exr:viii04-23}
Why is \(S^n\) noncontractible for \(n\ge1\)?
\end{exercise}
\begin{hint}
Use degree of the identity.

% VIII-PEDAGOGY-HINT-v1.1 exr:viii04-23 append
A contraction would make the identity homotopic to a constant sphere map. Compare their degrees, keeping the condition \(n\ge1\).
\end{hint}
\begin{solution}
A contraction would force degree \(1=0\), impossible.
\end{solution}

\begin{exercise}\label{exr:viii04-24}
What is the bridge to VIII/05?
\end{exercise}
\begin{hint}
Use projective filtration.

% VIII-PEDAGOGY-HINT-v1.1 exr:viii04-24 append
Normalize projective coordinates with nonzero last coordinate to \(1\). Identify this open \(\mathbb R^n\) chart as the new cell and the complementary projective subspace as the old skeleton.
\end{hint}
\begin{solution}
\(\mathbb RP^n\) is built from \(\mathbb RP^{n-1}\) by adding one \(n\)-cell.
\end{solution}


\section{Solved problem dossiers}

\begin{problem}[Antipodal parity]\label{prob:viii04-01}
Use degree to compare the antipodal map with the identity.
\end{problem}
\begin{solution}
For even \(n\), the degrees are \(-1\) and \(1\), so they are not homotopic.  For odd \(n\), both degrees are \(1\), so degree gives no obstruction.
\end{solution}

\begin{problem}[Circle antipodes]\label{prob:viii04-02}
Prove directly that every continuous \(f:S^1\to\mathbb R\) has equal values at some antipodal pair.
\end{problem}
\begin{solution}
Set \(h(x)=f(x)-f(-x)\).  On a semicircle from \(x\) to \(-x\), the endpoint values have opposite signs, so the intermediate value theorem gives \(h=0\).
\end{solution}

\begin{problem}[No odd map to \(S^0\)]\label{prob:viii04-03}
Prove there is no odd continuous map \(S^1\to S^0\).
\end{problem}
\begin{solution}
A continuous image of connected \(S^1\) in discrete \(S^0\) is a point, but a constant map cannot be odd.
\end{solution}

\begin{problem}[Equivalent Borsuk--Ulam form]\label{prob:viii04-04}
Derive the no-odd-map formulation and its converse.
\end{problem}
\begin{solution}
An odd \(S^n\to S^{n-1}\subset\mathbb R^n\) contradicts equal antipodal images. Conversely, an antipode-separating \(f\) yields the normalized odd difference map.
\end{solution}

\begin{problem}[Odd self-map]\label{prob:viii04-05}
Explain why an odd self-map of \(S^n\) cannot be null-homotopic.
\end{problem}
\begin{solution}
Its degree is odd and therefore nonzero, while a null-homotopic sphere map has degree zero.
\end{solution}

\begin{problem}[Antipodal \(S^2\)]\label{prob:viii04-06}
Compare the antipodal map on \(S^2\) with the identity and a constant map.
\end{problem}
\begin{solution}
Its degree is \(-1\), versus \(1\) for the identity and \(0\) for a constant; hence it is homotopic to neither.
\end{solution}

\begin{problem}[Antipodal \(S^3\)]\label{prob:viii04-07}
What does degree say about the antipodal map on \(S^3\)?
\end{problem}
\begin{solution}
Its degree is \(1\), equal to the identity degree.  Thus degree gives no obstruction to a homotopy with the identity; this chapter does not infer the homotopy from degree alone.
\end{solution}

\begin{problem}[No retraction]\label{prob:viii04-08}
Recover the no-retraction theorem from degree.
\end{problem}
\begin{solution}
If \(ri=\operatorname{id}_{S^n}\), then degree is \(1\); but \(i\) contracts through the disk, making \(ri\) null-homotopic of degree \(0\).
\end{solution}

\begin{problem}[Even-map quotient]\label{prob:viii04-09}
Prove the factorization of an even map through \(\mathbb RP^n\).
\end{problem}
\begin{solution}
Define \(\bar f([x])=f(x)\). Evenness gives well-definedness, and the quotient property of \(q\) gives continuity and uniqueness.
\end{solution}

\begin{problem}[Projective line]\label{prob:viii04-10}
Construct \(\mathbb RP^1\cong S^1\).
\end{problem}
\begin{solution}
Represent a line by angle \(\theta\) modulo \(\pi\) and send it to \(e^{2i\theta}\). This descends to a continuous bijection from compact \(\mathbb RP^1\) to Hausdorff \(S^1\).
\end{solution}

\begin{problem}[Quotient degree]\label{prob:viii04-11}
Compute the degree of \(S^1\to\mathbb RP^1\cong S^1\).
\end{problem}
\begin{solution}
It is \(z\mapsto z^2\), hence degree \(2\).
\end{solution}

\begin{problem}[Sphere from hemispheres]\label{prob:viii04-12}
Describe \(S^n\) as two disks glued along the boundary.
\end{problem}
\begin{solution}
Glue \(D^n_+\) and \(D^n_-\) by the identity on \(S^{n-1}\); the quotient is \(S^n\).
\end{solution}

\begin{problem}[One-point compactification]\label{prob:viii04-13}
Explain the stereographic model of \(S^n\).
\end{problem}
\begin{solution}
Stereographic projection identifies \(S^n\setminus\{N\}\) with \(\mathbb R^n\); the north pole becomes the point at infinity.
\end{solution}

\begin{problem}[Suspension sphere]\label{prob:viii04-14}
Explain geometrically why \(\Sigma S^n\cong S^{n+1}\).
\end{problem}
\begin{solution}
Use height in \([-1,1]\); at the two extreme heights all directions collapse to the suspension points.
\end{solution}

\begin{problem}[Antipodal orbit size]\label{prob:viii04-15}
Why is \(q:S^n\to\mathbb RP^n\) two-to-one?
\end{problem}
\begin{solution}
Each orbit is exactly \(\{x,-x\}\), and \(x
e -x\) on a unit sphere.
\end{solution}

\begin{problem}[Projective inclusion]\label{prob:viii04-16}
Explain \(\mathbb RP^{n-1}\subset\mathbb RP^n\).
\end{problem}
\begin{solution}
Lines lying in the coordinate hyperplane \(\mathbb R^n\subset\mathbb R^{n+1}\) form the embedded lower projective space.
\end{solution}

\begin{problem}[Projective open cell]\label{prob:viii04-17}
Why is \(\mathbb RP^n\setminus\mathbb RP^{n-1}\cong\mathbb R^n\)?
\end{problem}
\begin{solution}
Normalize a line not in the coordinate hyperplane so its last coordinate is \(1\); the remaining coordinates are arbitrary in \(\mathbb R^n\).
\end{solution}

\begin{problem}[Borsuk--Ulam application]\label{prob:viii04-18}
Apply Borsuk--Ulam to two continuous measurements on \(S^2\).
\end{problem}
\begin{solution}
For \(f=(f_1,f_2):S^2\to\mathbb R^2\), some antipodal pair has both measurements equal.
\end{solution}

\begin{problem}[Odd-vector obstruction]\label{prob:viii04-19}
Why can there be no continuous rule \(v:S^n\to S^{n-1}\) with \(v(-x)=-v(x)\)?
\end{problem}
\begin{solution}
That is exactly an odd map \(S^n\to S^{n-1}\), forbidden by Borsuk--Ulam.
\end{solution}

\begin{problem}[Bridge to cells]\label{prob:viii04-20}
Explain how projective geometry motivates cell attachment.
\end{problem}
\begin{solution}
The filtration \(\mathbb RP^{n-1}\subset\mathbb RP^n\) adds an open \(n\)-cell; VIII/05 abstracts the gluing of its closed disk along the boundary.
\end{solution}
